% Generated draft: requires mathematical and source review.
\documentclass[11pt]{article}
\usepackage[margin=25mm]{geometry}
\usepackage{fontspec}
\setmainfont{Latin Modern Roman}
\usepackage{amsmath,amssymb,mathtools,unicode-math}
\setmathfont{Latin Modern Math}
\usepackage{enumitem,needspace,xurl,hyperref}
\setlength{\parindent}{0pt}
\setlength{\parskip}{5pt}
\setlength{\emergencystretch}{3em}
\newcommand{\N}{\mathbb{N}}
\newcommand{\Z}{\mathbb{Z}}
\newcommand{\Q}{\mathbb{Q}}
\newcommand{\R}{\mathbb{R}}
\newcommand{\C}{\mathbb{C}}
\newcommand{\F}{\mathbb{F}}
\newcommand{\A}{\mathbb{A}}
\newcommand{\PP}{\mathbb{P}}
\newcommand{\E}{\mathbb{E}}
\newcommand{\eps}{\varepsilon}
\newcommand{\dd}{\,\mathrm{d}}
\begin{document}
\section*{UnsolvedMath ID: 30003625}
\textbf{Source problem code:} catalogue entry
\subsection{Definitions and mathematical statement}
The supplied question concerns perturbations of Jacobi matrices belonging to infinite-gap isospectral tori. It does not specify the operator domain, the definition of the torus, the precise Lieb--Thirring inequality, or the dependence of its constant. Consequently, the following gives a conditional conventional endpoint formulation; its conventions require confirmation from the cited report.

For this formulation, let \(E\subset\R\) be a nonempty compact set such that \((\min E,\max E)\setminus E\) has infinitely many connected components. A bounded two-sided Jacobi matrix is the self-adjoint operator on \(\ell^2(\Z)\) given by
\[
(Ju)_n=a_{n-1}u_{n-1}+b_nu_n+a_nu_{n+1},\qquad n\in\Z,
\]
where \(a_n>0\), \(b_n\in\R\), and both coefficient sequences are bounded. Let \(\mathcal T(E)\) denote a specified nonempty family of such operators with spectrum \(E\), intended to be the source's isospectral torus. The supplied record does not give the additional conditions defining that family; equality of spectra alone is not a definition of the intended torus.

For \(J_0\in\mathcal T(E)\), with coefficients \(a_n^{(0)},b_n^{(0)}\), and a Jacobi matrix \(J\) as above, define
\[
D_1(J,J_0)=\sum_{n\in\Z}\bigl(\lvert a_n-a_n^{(0)}\rvert+\lvert b_n-b_n^{(0)}\rvert\bigr).
\]
Write \(\sigma_{\mathrm d}(J)\) for the isolated eigenvalues of finite multiplicity and set \(\operatorname{dist}(x,E)=\inf_{t\in E}|x-t|\). One conventional endpoint question is whether
\[
\forall J_0\in\mathcal T(E)\ \exists C_{E,J_0}\in(0,\infty)\ \forall J\text{ with }D_1(J,J_0)<\infty:\qquad
\sum_{\lambda\in\sigma_{\mathrm d}(J)\setminus E}
\operatorname{dist}(\lambda,E)^{1/2}
\le C_{E,J_0}D_1(J,J_0),
\]
where eigenvalues are counted with multiplicity and the constant is independent of \(J\). A stronger version asks for one constant \(C_E\) valid for every \(J_0\in\mathcal T(E)\). The record does not distinguish these quantifiers. Calling this the endpoint \(p=1\) uses the convention that \(p\) indexes the coefficient perturbation summability and the spectral exponent is \(p-\tfrac12\); the record does not explicitly specify that convention.

The accompanying characterization problem asks which infinite-gap spectral sets admit the intended Lieb--Thirring bound. Its precise meaning remains unresolved until the allowed sets, the torus, the inequality, the exponent range, and the constant dependence are fixed. In particular, the two-sided domain and the displayed endpoint inequality are explicit conditional choices, not recovered source hypotheses.
\subsection{Short English statement}
The problem asks whether an endpoint Lieb--Thirring estimate controls eigenvalues created by perturbing Jacobi matrices in infinite-gap isospectral tori, and which spectral sets support such estimates. A conventional candidate controls the sum of square roots of eigenvalue distances by the total coefficient perturbation. The supplied record leaves essential conventions and quantifiers unresolved, so this formulation requires human mathematical review.
\subsection{Sources}
\begin{itemize}
\item Supplied local UnsolvedMath registry record, permanent ID \(30003625\), source problem code \texttt{catalogue entry}, title: Endpoint Lieb--Thirring Bounds for Infinite-Gap Jacobi Matrices. Collection: Oberwolfach Reports Open Problems. No individual proposer is supplied; the proposed year is recorded as 2017.
\item Supplied source citation: \emph{Mini-Workshop: Reflectionless Operators: The Deift and Simon Conjectures} (2018), Oberwolfach Reports, pp.\ 2943--2985, DOI \texttt{10.4171/owr/2017/49}. \url{https://doi.org/10.4171/owr/2017/49}. The literature-source list labels this reference Oberwolfach Report 49/2017. The supplied 2018 citation date and 2017 report designation are retained without reconciliation.
\item Supplied literature source: Christiansen--Zinchenko. \url{https://arxiv.org/abs/1609.09566}. No title, theorem number, or precise applicability statement is supplied.
\item Supplied literature assessment, checked on \texttt{2026-08-21}: bounds are described as known for periodic, finite-gap, and additional controlled infinite-gap settings, while the general infinite-gap endpoint problem and a full characterization of admissible sets are described as unresolved. This is supplied dated curation derived from the OWR record, not an independently verified literature assessment.
\item The registry labels the problem \texttt{open} and the statement \texttt{corrected\_verified}, explaining that notation, typography, or omitted definitions were repaired while preserving mathematical content. These are supplied registry labels, not independent verification of the formulation or current mathematical status. The supplied references have not been independently consulted for this draft.
\end{itemize}
\end{document}
